% Options for packages loaded elsewhere
\PassOptionsToPackage{unicode}{hyperref}
\PassOptionsToPackage{hyphens}{url}
\documentclass[
  polish,
  11pt,
  a4paper,
]{article}
\usepackage{xcolor}
\usepackage[margin=2cm]{geometry}
\usepackage{amsmath,amssymb}
\setcounter{secnumdepth}{-\maxdimen} % remove section numbering
\usepackage{iftex}
\ifPDFTeX
  \usepackage[T1]{fontenc}
  \usepackage[utf8]{inputenc}
  \usepackage{textcomp} % provide euro and other symbols
\else % if luatex or xetex
  \usepackage{unicode-math} % this also loads fontspec
  \defaultfontfeatures{Scale=MatchLowercase}
  \defaultfontfeatures[\rmfamily]{Ligatures=TeX,Scale=1}
\fi
\usepackage{lmodern}
\ifPDFTeX\else
  % xetex/luatex font selection
\fi
% Use upquote if available, for straight quotes in verbatim environments
\IfFileExists{upquote.sty}{\usepackage{upquote}}{}
\IfFileExists{microtype.sty}{% use microtype if available
  \usepackage[]{microtype}
  \UseMicrotypeSet[protrusion]{basicmath} % disable protrusion for tt fonts
}{}
\makeatletter
\@ifundefined{KOMAClassName}{% if non-KOMA class
  \IfFileExists{parskip.sty}{%
    \usepackage{parskip}
  }{% else
    \setlength{\parindent}{0pt}
    \setlength{\parskip}{6pt plus 2pt minus 1pt}}
}{% if KOMA class
  \KOMAoptions{parskip=half}}
\makeatother
\usepackage{longtable,booktabs,array}
\newcounter{none} % for unnumbered tables
\usepackage{calc} % for calculating minipage widths
% Correct order of tables after \paragraph or \subparagraph
\usepackage{etoolbox}
\makeatletter
\patchcmd\longtable{\par}{\if@noskipsec\mbox{}\fi\par}{}{}
\makeatother
% Allow footnotes in longtable head/foot
\IfFileExists{footnotehyper.sty}{\usepackage{footnotehyper}}{\usepackage{footnote}}
\makesavenoteenv{longtable}
\usepackage{graphicx}
\makeatletter
\newsavebox\pandoc@box
\newcommand*\pandocbounded[1]{% scales image to fit in text height/width
  \sbox\pandoc@box{#1}%
  \Gscale@div\@tempa{\textheight}{\dimexpr\ht\pandoc@box+\dp\pandoc@box\relax}%
  \Gscale@div\@tempb{\linewidth}{\wd\pandoc@box}%
  \ifdim\@tempb\p@<\@tempa\p@\let\@tempa\@tempb\fi% select the smaller of both
  \ifdim\@tempa\p@<\p@\scalebox{\@tempa}{\usebox\pandoc@box}%
  \else\usebox{\pandoc@box}%
  \fi%
}
% Set default figure placement to htbp
\def\fps@figure{htbp}
\makeatother
\ifLuaTeX
\usepackage[bidi=basic,shorthands=off]{babel}
\else
\usepackage[bidi=default,shorthands=off]{babel}
\fi
\ifLuaTeX
  \usepackage{selnolig} % disable illegal ligatures
\fi
\setlength{\emergencystretch}{3em} % prevent overfull lines
\providecommand{\tightlist}{%
  \setlength{\itemsep}{0pt}\setlength{\parskip}{0pt}}
\usepackage{fontspec}
\defaultfontfeatures{Ligatures=TeX}
\usepackage{amsmath,amssymb}
\usepackage{booktabs}
\usepackage{array}
\usepackage{geometry}
\usepackage{enumitem}
\usepackage{xcolor}
\usepackage{tcolorbox}
\usepackage{fancyhdr}
\usepackage{titlesec}
\usepackage{multicol}

\geometry{margin=1.8cm, top=2.5cm, bottom=2cm}
\pagestyle{fancy}
\fancyhf{}
\rhead{\small Projektowanie badań — 2026/27}
\lhead{\small Zarządzanie informacją | ISI UJ}
\cfoot{\thepage}

\titleformat{\section}{\large\bfseries\color{blue!70!black}}{\thesection.}{0.5em}{}
\titleformat{\subsection}{\normalsize\bfseries}{\thesubsection.}{0.5em}{}
\titleformat{\subsubsection}{\small\bfseries\color{gray!70!black}}{\thesubsubsection.}{0.5em}{}
\titlespacing*{\section}{0pt}{1.5ex}{1ex}
\titlespacing*{\subsection}{0pt}{1ex}{0.5ex}

\setlist[itemize]{nosep, leftmargin=1.5em}
\setlist[enumerate]{nosep, leftmargin=1.5em}

\tcbset{boxrule=0.5pt, left=4pt, right=4pt, top=4pt, bottom=4pt, fonttitle=\bfseries\small}

\newtcolorbox{definicja}[1][]{colback=blue!5, colframe=blue!50!black, title={#1}}
\newtcolorbox{uwagabox}[1][]{colback=orange!10, colframe=orange!70!black, title={#1}}
\newtcolorbox{wzorbox}[1][]{colback=gray!5, colframe=gray!50!black, title={#1}}
\newtcolorbox{przykladbox}[1][]{colback=purple!5, colframe=purple!50!black, title={#1}}
\usepackage{bookmark}
\IfFileExists{xurl.sty}{\usepackage{xurl}}{} % add URL line breaks if available
\urlstyle{same}
% fallback for those not using the hyperref driver hyperxmp:
\makeatletter
\@ifundefined{xmpquote}{\newcommand{\xmpquote}[1]{#1}}{}
\makeatother
\hypersetup{
  pdftitle={W04 --- Porównania: mapa projektu{,} efektu i testu},
  pdfauthor={Marek Deja},
  pdflang={pl-PL},
  hidelinks,
  pdfcreator={LaTeX via pandoc}}

\title{W04 --- Porównania: mapa projektu, efektu i testu}
\usepackage{etoolbox}
\makeatletter
\providecommand{\subtitle}[1]{% add subtitle to \maketitle
  \apptocmd{\@title}{\par {\large #1 \par}}{}{}
}
\makeatother
\subtitle{Zarządzanie informacją --- handout --- 2026/27}
\author{Marek Deja}
\date{2026/27}

\begin{document}
\maketitle

\section{Najpierw ustal strukturę}\label{najpierw-ustal-strukturux119}

{\def\LTcaptype{none} % do not increment counter
\begin{longtable}[]{@{}
  >{\raggedright\arraybackslash}p{(\linewidth - 4\tabcolsep) * \real{0.3333}}
  >{\raggedright\arraybackslash}p{(\linewidth - 4\tabcolsep) * \real{0.3333}}
  >{\raggedright\arraybackslash}p{(\linewidth - 4\tabcolsep) * \real{0.3333}}@{}}
\toprule\noalign{}
\begin{minipage}[b]{\linewidth}\raggedright
Projekt
\end{minipage} & \begin{minipage}[b]{\linewidth}\raggedright
Wielkość główna
\end{minipage} & \begin{minipage}[b]{\linewidth}\raggedright
Kluczowa kontrola
\end{minipage} \\
\midrule\noalign{}
\endhead
\bottomrule\noalign{}
\endlastfoot
Dwie niezależne grupy, wynik ilościowy & Różnica średnich & Inne osoby,
N i SD w obu grupach \\
Dwa pomiary tych samych osób & Średnia różnica wewnątrz par & ID i
kompletne pary \\
Dwie zmienne kategorialne & Proporcje i związek w tabeli & Liczebności,
mianowniki, oczekiwania \\
Więcej grup & Zaplanowany kontrast lub porównanie ogólne & Wielokrotność
i właściwy model \\
Wynik interpretowany rangowo & Położenie/rangi według metody & Nie
utożsamiać z testem średnich \\
\end{longtable}
}

Równe liczebności nie tworzą par. Wielokrotne zadania tej samej osoby
nie są niezależnymi osobami. Grupy obserwacyjne nie stają się
randomizowane przez użycie permutacji.

\section{Jeden kierunek we wszystkich
wynikach}\label{jeden-kierunek-we-wszystkich-wynikach}

\[\Delta=\mu_N-\mu_D,\qquad
\widehat\Delta=\bar{x}_N-\bar{x}_D.\]

Nowi minus doświadczeni: dodatni wynik czasu oznacza dłuższy czas
nowych, dodatni indeks --- wyższą deklarację zgodnie z kierunkiem skali.
Po odwróceniu kolejności zmieniamy znak różnicy, t, g i granice CI:
{[}L; U{]} przechodzi w {[}−U; −L{]}. Dwustronne p pozostaje takie samo.

\section{Welch: obie średnie mają
niepewność}\label{welch-obie-ux15brednie-majux105-niepewnoux15bux107}

\[
SE=\sqrt{\frac{s_N^2}{n_N}+\frac{s_D^2}{n_D}},
\quad t_W=\frac{\widehat\Delta}{SE},
\quad CI_{95\%}=\widehat\Delta\pm t_{0{,}975;\nu}SE.
\]

Welch nie wymaga równych wariancji, ale nadal wymaga właściwego projektu
i uzasadnienia przybliżenia. Niecałkowite df jest normalnym wynikiem
przybliżenia, nie błędem liczebności. Test normalności nie ocenia
niezależności i nie jest uniwersalnym przełącznikiem metod.

Przykład hipotetyczny: po 5 osób, średnie czasu 12 i 8 minut, SD w obu
grupach około 3,162. Różnica wynosi 4 minuty, SE=2, t(8)=2, p około
0,081, CI {[}−0,61; 8,61{]} minuty. To nieprecyzyjna dodatnia estymata,
nie dowód równości.

\[
s_p=\sqrt{\frac{(n_N-1)s_N^2+(n_D-1)s_D^2}{n_N+n_D-2}},
\quad
g\approx\left(1-\frac{3}{4(n_N+n_D)-9}\right)
\frac{\widehat\Delta}{s_p}.
\]

G dzieli różnicę przez rozrzut osób, nie przez SE. Naturalna różnica w
minutach lub punktach pozostaje podstawą sensu praktycznego. Duże g i p
powyżej 0,05 mogą wystąpić jednocześnie przy małej próbie.

\section{Pary: jedna różnica na
osobę}\label{pary-jedna-ruxf3ux17cnica-na-osobux119}

\[
d_i=x_{i,\mathrm{po}}-x_{i,\mathrm{przed}},
\quad t=\frac{\bar d}{s_d/\sqrt{n_{\mathrm{par}}}},
\quad df=n_{\mathrm{par}}-1.
\]

Ujemna zmiana czasu oznacza skrócenie. Założenie normalności dla
dokładnego modelu dotyczy różnic, nie osobno kolumn przed i po. Bez
kontroli zmiana może odzwierciedlać także uczenie się, a nie tylko
interwencję.

\section{Resampling musi pasować do
projektu}\label{resampling-musi-pasowaux107-do-projektu}

Permutacja dwóch grup przestawia etykiety przy zachowaniu liczebności,
jeśli pod H0 dopuszcza to wymienialność lub mechanizm randomizacji. Sama
równość średnich nie wystarcza przy dowolnie różnych rozkładach. Surowe
tasowanie różnicy nie dziedziczy automatycznie odporności Welcha na
nierówne wariancje.

Bootstrap dwóch grup losuje z powtórzeniami osobno w każdej grupie, aby
szacować niepewność różnicy. Dla par zachowujemy całą parę. Zmiana
znaków różnic ma własne uzasadnienie symetrii lub randomizacji; nie jest
dowolnym mieszaniem wszystkich wartości.

\section{Kategorie: od komórki do
statystyki}\label{kategorie-od-komuxf3rki-do-statystyki}

\[E_{ij}=\frac{n_{i\cdot}n_{\cdot j}}N,\qquad
\chi^2=\sum_{i,j}\frac{(O_{ij}-E_{ij})^2}{E_{ij}},
\qquad df=(r-1)(c-1).\]

O to liczebność obserwowana, E oczekiwana przy niezależności. E może być
ułamkiem. Do testu podajemy liczebności, nie procenty. Małe oczekiwania
mogą pogarszać przybliżenie chi-kwadrat; Fisher lub odpowiednie Monte
Carlo wymagają jawnego opisu procedury.

Monte Carlo w omawianym teście tabel zachowuje marginesy. Fisher 2 na 2
ma interpretację warunkową przy tych marginesach. Słowo „dokładny'' nie
oznacza braku wymagań dotyczących projektu lub niezależności.

\section{Wielkość związku nie jest
p}\label{wielkoux15bux107-zwiux105zku-nie-jest-p}

\[V=\sqrt{\frac{\chi^2}{N\min(r-1,c-1)}}.\]

V jest nieujemne i nie ma kierunku. Nie jest procentem skuteczności.
Prosty percentylowy bootstrap V może nie obejmować zera także przy
niezależności populacyjnej, dlatego nie używamy jego dolnej granicy jako
automatycznego testu związku.

W przykładzie 18/30 wobec 32/40:

{\def\LTcaptype{none} % do not increment counter
\begin{longtable}[]{@{}lrl@{}}
\toprule\noalign{}
Miara: nowi względem doświadczonych & Wynik & Odczyt \\
\midrule\noalign{}
\endhead
\bottomrule\noalign{}
\endlastfoot
Różnica proporcji & −0,20 & O 20 punktów procentowych mniej \\
Iloraz proporcji & 0,75 & Udział stanowi 75\% udziału drugiej grupy \\
Iloraz szans & 0,375 & Inny iloraz: sukcesy do niepowodzeń \\
\end{longtable}
}

Przedział każdej miary trzeba nazwać. Wartością zerową różnicy jest 0, a
brakiem związku dla ilorazów jest 1. Nie przenosimy granic między
skalami bez odpowiedniego przekształcenia.

\section{Rangi i wiele grup: granice
skrótu}\label{rangi-i-wiele-grup-granice-skruxf3tu}

Test Manna--Whitneya nie jest automatycznym testem median. Interpretacja
przesunięcia położenia wymaga dodatkowych warunków. Porównanie rang może
odpowiadać innemu pytaniu niż przeciętny koszt czasu.

ANOVA pyta, czy wszystkie średnie są równe; wynik ogólny nie wskazuje
pary. Zaplanowane kontrasty i analizy po obejrzeniu wyników wymagają
odpowiednich reguł wielokrotności. Istotny wynik ogólny nie uprawnia do
dowolnego wybierania kolejnych p.

W raporcie: projekt, grupy, N, wielkość i kierunek efektu, jednostka,
CI, statystyka, df, p i ograniczenie. Wniosek o przyczynie wymaga
uzasadnienia projektu, nie tylko małego p.

Definicje efektów:
\href{https://doi.org/10.3389/fpsyg.2013.00863}{Lakens (2013)}. Warunki
przestawień:
\href{https://pmc.ncbi.nlm.nih.gov/articles/PMC4010955/}{Winkler i in.
(2014)}.

\end{document}
